\begin{remark}[E18: the whole curve, against the spans the corpus declares]%
\label{rem:rollup-spans}
E4 read the automatic interval at five query ranges and
Remark~\ref{rem:rollup-default-cost} argues from them. Five points are enough to say the
interval is a function of the range and not enough to say which function --- and
\S\ref{sec:threats} called the emitter's assumed $60$\,s ``systematic rather than
arbitrary'' without asking how often it is right. It is answerable, because a Datadog
widget records the span it is viewed over. The corpus's dashboards declare $460$ of
them, in ten distinct values, and $60$\,s is the interval of exactly one:
four hours, which is $28.3\%$ of them. The modal span is one hour at $38.3\%$, where the
interval is $20$\,s and every translated rollup is therefore out by a factor of three.

So all ten were measured, on one series of an hour's points at $10$\,s --- the interval
depends on the range and not on the data, so an hour of backfill is enough to ask about
a month. E4's five are re-asked as the control rather than taken on trust.

\begin{center}\footnotesize
\begin{tabular}{@{}lrrrrrrrrrr@{}}
\toprule
Range & \texttt{1m} & \texttt{5m} & \texttt{10m} & \texttt{15m} & \texttt{30m}
      & \texttt{1h} & \texttt{4h} & \texttt{1d} & \texttt{1w} & \texttt{1mo} \\
\midrule
Reported (s) & $1$ & $2$ & $2$ & $5$ & $10$ & $20$ & $60$ & $300$ & $3{,}600$
             & $14{,}400$ \\
Observed (s) & $10$ & $10$ & $10$ & $10$ & $10$ & $20$ & $60$ & $300$ & $3{,}600$
             & --- \\
\bottomrule
\end{tabular}
\end{center}

\noindent
The curve is monotone and all five of E4's values reappear, so nothing in
Remark~\ref{rem:rollup-default-cost} moves and the two runs read as one table. The
second row is the part a translator has to care about, and it was not new: E4's own
cache recorded the same divergence at two of its five ranges and nobody read it, which
is the more useful lesson in this remark. An experiment answers the question it was
designed around; the rest of what it captured sits there until someone asks.

The interval a response \emph{reports} is not always the interval the series \emph{has}.
These points were submitted $10$\,s apart, and at every range below an hour the reported
interval is finer than that while the returned points stay $10$\,s apart: the field is
the granularity the service resolved the request to, not a description of what came
back. Recovering $\Delta$ from the response rather than from the dashboard is the
obvious way to avoid the assumption this remark is about, and it does not work --- at a
one-minute range it reads $1$ where the data is at $10$, and divides by it. What
\texttt{ROLLUP\_DEFAULT} marks as missing from the query is missing from the answer too.

The month column is blank because this run's series is $3{,}500$\,s long, so a
month-wide window holds one bucket of it and no spacing can be observed. The reported
$14{,}400$ stands unchecked against data, which is the honest state for that one cell
and the reason the row is dashed rather than filled
(\texttt{scripts\slash trial\_checks.py e18}). It is also the cell the emitter's
$60$\,s misses by the widest margin --- a $240$th of the reported interval, against a
sixtieth at one week --- so the largest error this table implies is the one value in it
with nothing observed behind it.
\end{remark}
